%% ==========================================================================
%% smpost/profile_photo_youtube.tex : the JitJet YouTube profile picture, drawn
%% with TikZ from the same cone as cover/cover_timeline.tex.
%%
%% Compile and rasterise:   make profile-youtube
%%   (or: cd smpost && pdflatex profile_photo_youtube
%%        gs -q -dBATCH -dNOPAUSE -sDEVICE=png16m -r254 -g800x800 -dPDFFitPage \
%%           -dTextAlphaBits=4 -dGraphicsAlphaBits=4 \
%%           -sOutputFile=profile_photo_youtube.png profile_photo_youtube.pdf)
%% then upload smpost/profile_photo_youtube.png as the channel picture.
%%
%% YouTube wants a square of at least 800 x 800 px and shows it as a circle,
%% mostly very small (98 px next to a video, 160 px on the channel page). The
%% page is 8 cm x 8 cm, 1 px = 0.01 cm, and everything that matters sits inside
%% the inscribed circle of radius 4 cm around the centre; the corners are cut.
%%
%% Layout: a gold rim just inside the crop; the wordmark JitJet across the
%% upper half, as wide as the circle allows, with the subtitle under it; the
%% cone of the cover, horizontal across the lower half as the emblem (glow,
%% b quark and shower, rings, the R = 0.4 jet, the b stamp, the three-prong
%% top cap), stripped of the fine detail that would not survive a thumbnail.
%% ==========================================================================
\documentclass[tikz,border=0pt]{standalone}
\usepackage[T1]{fontenc}
\usepackage{lmodern}
\usepackage{tgheros}          % TeX Gyre Heros for the sans-serif wordmark
\usepackage{amssymb}
\usetikzlibrary{calc,arrows.meta,decorations.pathmorphing}

%% -------- palette (identical to cover.tex) ----------------------------------
\definecolor{navy}{HTML}{0B1D3A}
\definecolor{navytop}{HTML}{17365F}
\definecolor{gold}{HTML}{F4B942}
\definecolor{bred}{HTML}{FF5A5A}
\definecolor{bblue}{HTML}{6E8EFF}
\definecolor{fsr}{HTML}{3DD16E}
\definecolor{isr}{HTML}{FF9A3C}
\definecolor{hadroncol}{HTML}{FFC24D}
\definecolor{trackcol}{HTML}{E4ECF7}
\definecolor{ecalcol}{HTML}{3FBF77}
\definecolor{hcalcol}{HTML}{E58F2A}
\definecolor{muoncol}{HTML}{8FB8FF}

%% -------- the cone (as in cover_timeline.tex) --------------------------------
\def\kk{0.30}
\pgfmathdeclarefunction{rf}{1}{\pgfmathparse{0.18*#1+0.6*(1+((#1-16.8)/1.2)/sqrt(1+((#1-16.8)/1.2)^2))}}
\pgfmathdeclarefunction{sc}{1}{\pgfmathparse{min(rf(#1)/2.2,1)}}
\newcommand{\xy}[3]{({#1+\kk*#2*sc(#1)*cos(#3)},{#2*sc(#1)*sin(#3)})}
\newcommand{\xring}[2]{%
  \draw[ringback] ({#1},{#2}) arc[start angle=90, end angle=-90, x radius={\kk*#2}, y radius={#2}];
  \draw[ring]     ({#1},{#2}) arc[start angle=90, end angle=270, x radius={\kk*#2}, y radius={#2}];}
\newcommand{\xellipse}[3][]{\path[#1] ({#2},0) ellipse [x radius={\kk*#3}, y radius={#3}];}

%% where the cone sits (page centre is (4,4)) and its size: 1 unit of s = \cs cm
\def\cx{1.1}
\def\cy{2.55}
\def\cs{0.26}

\tikzset{
  ring/.style={draw=gold!75!white, line width=0.6pt, opacity=0.9},
  ringback/.style={draw=white, line width=0.4pt, opacity=0.3, dash pattern=on 1pt off 1pt},
  meridian/.style={draw=white, line width=0.3pt, opacity=0.15},
  edge/.style={draw=gold!80!white, line width=1.1pt},
  bline/.style={draw=bred, line width=1.8pt, line cap=round},
  bbarline/.style={draw=bblue, line width=1.8pt, line cap=round},
  gluon/.style={draw=fsr, line width=0.9pt, decorate,
    decoration={coil, aspect=0.55, segment length=2.2pt, amplitude=1.3pt}},
  isrgluon/.style={gluon, draw=isr},
  quark/.style={draw=hadroncol, line width=0.9pt, line cap=round},
  track/.style={draw=trackcol, line width=0.6pt, opacity=0.9},
  pf/.style={circle, inner sep=0pt, minimum size=3.6pt, draw=navy, line width=0.25pt},
  roadline/.style={draw=gold, line width=0.9pt, dash pattern=on 4pt off 3pt},
}

\begin{document}
\begin{tikzpicture}[x=1cm,y=1cm]
\useasboundingbox (0,0) rectangle (8,8);
\clip (0,0) rectangle (8,8);

%% ======== background and rim ================================================
\shade[inner color=navytop, outer color=navy!80!black] (0,0) rectangle (8,8);
\draw[white, opacity=0.05, line width=0.3pt, step=0.5cm] (0,0) grid (8,8);
\draw[gold, line width=2.2pt] (4,4) circle[radius=3.82];
\draw[gold!60!navy, line width=0.6pt] (4,4) circle[radius=3.62];

%% ======== the cone ==========================================================
\begin{scope}[shift={(\cx,\cy)}, scale=\cs]

\shade[inner color=gold, outer color=navy, opacity=0.75] (0,0) circle[radius=1.6];

\fill[gold, opacity=0.13]
  plot[domain=0:19, samples=80] (\x,{rf(\x)})
  arc[start angle=90, end angle=-90, x radius={\kk*rf(19)}, y radius={rf(19)}]
  -- plot[domain=19:0, samples=80] (\x,{-rf(\x)}) -- cycle;
\foreach \t in {0,45,...,315}
  \draw[meridian] plot[domain=0:19, samples=70] ({\x+\kk*rf(\x)*cos(\t)},{rf(\x)*sin(\t)});
\draw[edge] plot[domain=0:19, samples=80] (\x,{rf(\x)});
\draw[edge] plot[domain=0:19, samples=80] (\x,{-rf(\x)});
\draw[roadline] (0,0) -- (19,0);

%% the family is born: the sibling leaves, the b showers along the axis
\draw[bbarline] (0,0) -- (0.5,2.2);
\draw[isrgluon] (0,0) -- (1.9,1.9);
\coordinate (B1) at (1.3,0.12);
\coordinate (B2) at (2.8,0.02);
\coordinate (B3) at (4.6,-0.08);
\path \xy{3.4}{1.0}{120} coordinate (G1);
\path \xy{4.4}{1.05}{210} coordinate (G2);
\path \xy{4.4}{1.1}{330} coordinate (G3);
\path \xy{4.4}{1.5}{112} coordinate (Q1);
\path \xy{4.4}{1.0}{145} coordinate (Q2);
\draw[bline] (0,0) -- (B1) -- (B2) -- (B3);
\draw[gluon] (B1) -- (G1);
\draw[gluon] (B2) -- (G2);
\draw[gluon] (B2) -- (G3);
\draw[quark] (G1) -- (Q1);
\draw[quark] (G1) -- (Q2);
\xring{4.4}{rf(4.4)}

%% the hadrons fly on, parallel to the axis
\foreach \r/\t in {1.4/118, 0.9/150, 1.2/335, 1.6/95, 0.5/30, 1.0/200, 0.6/305}
  \draw[track] \xy{6.5}{\r}{\t} -- \xy{12}{\r}{\t};
\xring{8}{rf(8)}

%% the census and the family inside R = 0.4
\foreach \r/\t/\c in {1.4/118/trackcol, 0.9/150/trackcol, 1.2/335/trackcol, 1.6/95/trackcol,
                      0.5/30/trackcol, 1.0/200/ecalcol, 1.4/222/ecalcol, 0.6/305/hcalcol, 0.5/250/muoncol}
  \node[pf, fill=\c] at \xy{12}{\r}{\t} {};
\xring{12}{rf(12)}
\xellipse[fill=gold, fill opacity=0.15, draw=gold, line width=0.9pt]{14.5}{1.75}
\foreach \r/\t/\c in {1.4/118/trackcol, 0.9/150/trackcol, 1.2/335/trackcol, 1.6/95/trackcol,
                      0.5/30/trackcol, 1.0/200/ecalcol, 1.4/222/ecalcol, 0.6/305/hcalcol, 0.5/250/muoncol}
  \node[pf, fill=\c] at \xy{14.5}{\r}{\t} {};
\xring{14.5}{rf(14.5)}

%% the top family: the AK8 cap of the cone
\coordinate (T) at (19,0);
\xellipse[fill=gold, fill opacity=0.2, draw=gold, line width=0.9pt, dash pattern=on 2.5pt off 2pt]{19}{rf(19)}
\foreach \t/\c in {90/bred, 210/hadroncol, 330/hadroncol}{
  \path \xy{19}{2.25}{\t} coordinate (SJ);
  \draw[draw=\c, fill=\c, fill opacity=0.2, line width=0.8pt] (SJ) ellipse [x radius={\kk*1.3}, y radius=1.3];
  \draw[draw=\c, line width=0.7pt, opacity=0.85] (T) -- (SJ);
  \pgfmathsetseed{\t}
  \foreach \i in {1,...,6}{
    \pgfmathsetmacro{\rr}{1.0*sqrt(rnd)}
    \pgfmathsetmacro{\tt}{360*rnd}
    \node[pf, fill=\c, minimum size=3pt] at ({19+\kk*(2.25*cos(\t)+\rr*cos(\tt))},{2.25*sin(\t)+\rr*sin(\tt)}) {};}}

%% the two stamps on the axis
\node[circle, draw=gold, fill=bred, line width=1pt, minimum size=13pt, inner sep=0pt,
      font=\sffamily\bfseries\scriptsize, text=white] at (16.6,0) {b};
\node[circle, fill=gold, draw=navy, line width=0.6pt, minimum size=13pt, inner sep=0pt,
      font=\sffamily\bfseries\scriptsize, text=navy] at (T) {$t$};

\end{scope}

%% ======== the wordmark ======================================================
\node[anchor=center] at (4,5.5)
  {\fontsize{60}{64}\selectfont\sffamily\bfseries{\color{white}Jit}{\color{gold}Jet}};
\draw[gold, line width=0.7pt] (2.2,4.65) -- (5.8,4.65);
\node[anchor=center, white, font=\fontsize{20}{13}\selectfont\sffamily] at (4,4.25) {Journey Is The Jet};
\node[anchor=center, gold!90!white, font=\fontsize{12}{10}\selectfont\sffamily\itshape] at (4,3.70) {(bottom to top)};

\end{tikzpicture}
\end{document}
